\chapter{Branching and Merging}\label{ch:git-branch}

A branch is a file containing a hash (\autoref{sec:git-refs}). Everything in
this chapter follows from that, including the parts that frighten people.

\section{Creating and Switching}\label{sec:git-branch}

\begin{keys}[0.34]
	\cmd{git branch}                  & List local branches; \texttt{*} marks the current \\
	\cmd{git branch -a}               & Include remote-tracking branches                  \\
	\cmd{git branch -v}               & \dots{} with each one's latest commit             \\
	\cmd{git switch feature}          & Move to an existing branch                        \\
	\cmd{git switch -c feature}       & Create it and move there                          \\
	\cmd{git switch -c fix main}      & \dots{} starting from \cmd{main} rather than here \\
	\cmd{git switch -}                & Back to the previous branch                       \\
	\cmd{git branch -d feature}       & Delete it --- refuses if unmerged                 \\
	\cmd{git branch -D feature}       & Delete it anyway                                  \\
	\cmd{git branch -m old new}       & Rename                                            \\
\end{keys}

\begin{note}
	\cmd{git switch} and \cmd{git restore} were introduced in Git 2.23 to split the
	two unrelated jobs that \cmd{git checkout} had been doing --- moving between
	branches, and throwing away file changes. \cmd{git checkout} still works and
	you will see it everywhere; the newer pair is clearer, and much harder to use
	by accident.
\end{note}

\begin{eg}
	Branching costs nothing, so branch for everything:
\begin{shell}
$ git switch -c fix-config-crash
Switched to a new branch 'fix-config-crash'
$ # ... edit, commit, edit, commit ...
$ git switch main
$ git merge fix-config-crash
$ git branch -d fix-config-crash
\end{shell}
\end{eg}

\section{Fast-Forward and Three-Way Merges}\label{sec:git-merge}

\begin{definition}[Fast-forward]
	When the target branch is a direct ancestor of the one being merged, there is
	nothing to combine: Git simply moves the pointer forward. No merge commit is
	created.
\end{definition}

\begin{center}
	\begin{tikzpicture}[
			x=1mm, y=1mm, font=\footnotesize\ttfamily,
			c/.style={draw=Gray!60, circle, inner sep=1.5pt, fill=Gray!5, minimum size=6.5mm},
			ref/.style={draw=SeaGreen!60!black, rounded corners=2pt, inner sep=2pt,
					fill=SeaGreen!8, font=\scriptsize\ttfamily},
			e/.style={-{Stealth[length=4pt]}, Gray!70!black}
		]
		% before
		\node[font=\scriptsize\sffamily, Gray!60!black] at (-14, 4) {before};
		\node[c] (a) at (  0,  4) {A};
		\node[c] (b) at ( 18,  4) {B};
		\node[c] (c) at ( 36,  4) {C};
		\draw[e] (b) -- (a); \draw[e] (c) -- (b);
		\node[ref] at (  0, 16) {main};
		\node[ref] at ( 36, 16) {feature};
		% after
		\node[font=\scriptsize\sffamily, Gray!60!black] at (-14, -22) {after};
		\node[c] (a2) at (  0, -22) {A};
		\node[c] (b2) at ( 18, -22) {B};
		\node[c] (c2) at ( 36, -22) {C};
		\draw[e] (b2) -- (a2); \draw[e] (c2) -- (b2);
		\node[ref] at ( 36, -10) {main, feature};
	\end{tikzpicture}
\end{center}

When both branches have moved on, Git makes a \vocab{merge commit} with two
parents:

\begin{center}
	\begin{tikzpicture}[
			x=1mm, y=1mm, font=\footnotesize\ttfamily,
			c/.style={draw=Gray!60, circle, inner sep=1.5pt, fill=Gray!5, minimum size=6.5mm},
			m/.style={draw=MidnightBlue!60, circle, inner sep=1.5pt,
					fill=MidnightBlue!8, minimum size=6.5mm},
			ref/.style={draw=SeaGreen!60!black, rounded corners=2pt, inner sep=2pt,
					fill=SeaGreen!8, font=\scriptsize\ttfamily},
			e/.style={-{Stealth[length=4pt]}, Gray!70!black}
		]
		\node[c] (a) at (  0,   0) {A};
		\node[c] (b) at ( 18,   0) {B};
		\node[c] (d) at ( 40,  10) {D};
		\node[c] (c) at ( 40, -10) {C};
		\node[m] (m) at ( 62,   0) {M};
		\draw[e] (b) -- (a);
		\draw[e] (d) -- (b);
		\draw[e] (c) -- (b);
		\draw[e] (m) -- (d);
		\draw[e] (m) -- (c);
		\node[ref] at ( 62,  12) {main};
		\node[ref] at ( 40, -21) {feature};
	\end{tikzpicture}
\end{center}

\begin{keys}[0.34]
	\cmd{git merge feature}          & Merge \cmd{feature} into the current branch \\
	\cmd{git merge -{}-no-ff feature} & Always make a merge commit, even if it could fast-forward \\
	\cmd{git merge -{}-squash feature} & Take the changes as one uncommitted lump   \\
	\cmd{git merge -{}-abort}         & Give up and return to where you started    \\
\end{keys}

\begin{note}
	\cmd{--no-ff} is worth knowing. A fast-forward loses the fact that a group of
	commits belonged to one branch; a merge commit records it, so
	\cmd{git log --graph} keeps showing the shape of the work. Many teams require
	it, and GitHub's ``Create a merge commit'' button does exactly this.
\end{note}

\section{Conflicts}\label{sec:git-conflict}

\begin{definition}[Conflict]
	Both sides changed the same region of the same file, and Git will not guess.
	It writes both versions into the file with markers and stops, leaving you to
	produce the correct result.
\end{definition}

\begin{conf}
<<<<<<< HEAD
    timeout = 30;
=======
    timeout = 60;
>>>>>>> feature
\end{conf}

\begin{keys}[0.30]
	\texttt{<{}<{}<{}<{}<{}<{}< HEAD}   & Everything below is \emph{your} current branch   \\
	\texttt{=======}        & The divider                                      \\
	\texttt{>{}>{}>{}>{}>{}>{}> feature} & Everything above is the branch being merged      \\
\end{keys}

The procedure is always the same:

\begin{enumerate}[(1)]
	\item \cmd{git status} lists the conflicted files.
	\item Open each one, delete the markers, and leave the text you actually want
	      --- which may be either side, both, or something new.
	\item \cmd{git add} the file to say it is resolved.
	\item \cmd{git commit} when they are all done, or \cmd{git merge --abort} to
	      walk away.
\end{enumerate}

\begin{keys}[0.32]
	\cmd{git status}                 & Which files are still conflicted              \\
	\cmd{git diff}                   & The remaining conflicts, in a combined diff   \\
	\cmd{git checkout -{}-ours file} & Take your side wholesale                      \\
	\cmd{git checkout -{}-theirs file} & Take theirs wholesale                       \\
	\cmd{git merge -{}-abort}        & Undo the whole merge                          \\
	\cmd{git mergetool}              & Open a three-way merge tool --- \cmd{nvim -d} by default with \cmd{merge.tool=vimdiff} \\
\end{keys}

\begin{note}
	``Ours'' and ``theirs'' mean the opposite of what you expect during a
	\emph{rebase}: rebasing replays your commits on top of the other branch, so
	\cmd{--ours} is the branch you are rebasing onto. Check with
	\cmd{git status}, which names both sides explicitly, rather than trusting the
	words.
\end{note}

\begin{moral}
	A conflict is not an error and not a failure. It is Git declining to guess,
	which is the correct behaviour. The way to have fewer of them is to merge
	\emph{more} often, not less: two branches that diverge for a month conflict
	badly, and two that sync daily barely conflict at all.
\end{moral}

\section{Rebase}\label{sec:git-rebase}

\begin{definition}[Rebase]
	Replay a series of commits onto a different base. Git takes each commit on
	your branch, computes its diff, and applies it to the new starting point ---
	producing \emph{new} commits with new hashes.
\end{definition}

\begin{center}
	\begin{tikzpicture}[
			x=1mm, y=1mm, font=\footnotesize\ttfamily,
			c/.style={draw=Gray!60, circle, inner sep=1.5pt, fill=Gray!5, minimum size=6.5mm},
			n/.style={draw=WildStrawberry!60, circle, inner sep=1.5pt,
					fill=WildStrawberry!8, minimum size=6.5mm},
			ref/.style={draw=SeaGreen!60!black, rounded corners=2pt, inner sep=2pt,
					fill=SeaGreen!8, font=\scriptsize\ttfamily},
			e/.style={-{Stealth[length=4pt]}, Gray!70!black}
		]
		\node[c] (a) at (  0,  0) {A};
		\node[c] (b) at ( 18,  0) {B};
		\node[c] (d) at ( 36,  0) {D};
		\node[n] (c) at ( 54,  0) {C'};
		\draw[e] (b) -- (a); \draw[e] (d) -- (b); \draw[e] (c) -- (d);
		\node[ref] at ( 36, 12) {main};
		\node[ref] at ( 54, -12) {feature};
		\node[font=\scriptsize\sffamily, Gray!60!black, align=left]
		at (95, 0) {C became C'\\ --- same change,\\ new commit};
	\end{tikzpicture}
\end{center}

\begin{keys}[0.34]
	\cmd{git rebase main}            & Replay this branch on top of \cmd{main}      \\
	\cmd{git rebase -{}-continue}    & After resolving a conflict                   \\
	\cmd{git rebase -{}-skip}        & Drop the commit being applied                \\
	\cmd{git rebase -{}-abort}       & Return to where you started                  \\
	\cmd{git pull -{}-rebase}        & Fetch, then rebase rather than merge         \\
\end{keys}

\begin{intuition}
	Merge preserves what happened; rebase produces what you wish had happened. A
	merge records that two lines of work existed and came together. A rebase makes
	it look as though you had started from the latest \cmd{main} all along, giving
	a linear history that is far easier to read and to bisect.
\end{intuition}

\begin{moral}[The golden rule]
	\textbf{Never rebase commits that other people have.} Rebasing creates new
	commits with new hashes; anyone who had the old ones now has a divergent
	history, and merging the two duplicates every commit. Rebase your own
	unpublished work freely; once it is pushed to a shared branch, merge instead.
\end{moral}

\section{Interactive Rebase}\label{sec:git-irebase}

\cmd{git rebase -i} is where a messy series of commits becomes a clean one.

\begin{shell}
$ git rebase -i main
\end{shell}

Your editor opens with the commits to be replayed, oldest first:

\begin{conf}
pick 9c4e7d2 Extract parse_line
pick 3f8a2b1 fix typo
pick 7b1e4a9 Add config validation
pick 2d8c5f0 oops, forgot the header

# Commands:
#  p, pick   = use commit
#  r, reword = use commit, but edit the message
#  e, edit   = use commit, but stop for amending
#  s, squash = meld into previous commit, keep both messages
#  f, fixup  = meld into previous commit, discard this message
#  d, drop   = remove commit
#  (lines may be reordered; they are applied top to bottom)
\end{conf}

Editing it into this:

\begin{conf}
pick 9c4e7d2 Extract parse_line
fixup 3f8a2b1 fix typo
pick 7b1e4a9 Add config validation
fixup 2d8c5f0 oops, forgot the header
\end{conf}

turns four commits into two clean ones, each self-contained.

\begin{keys}[0.30]
	\cmd{reword}     & Fix a message without touching the content                  \\
	\cmd{squash}     & Combine, and edit the joined message                         \\
	\cmd{fixup}      & Combine, discarding this commit's message. The common case   \\
	\cmd{drop}       & Remove a commit entirely                                     \\
	\cmd{edit}       & Stop there so you can amend or split it                      \\
	Reorder lines    & Reorder the commits                                          \\
\end{keys}

\begin{note}
	\cmd{git commit --fixup=<sha>} marks a commit as a fix for an earlier one, and
	\cmd{git rebase -i --autosquash main} then pre-arranges the list for you. Set
	\cmd{git config --global rebase.autosquash true} and the workflow becomes: fix
	something, mark it, and squash the whole series before opening a pull request.
\end{note}

\begin{moral}
	Commit early and often while working --- including ``wip'' and ``oops''. Then
	tidy the series with an interactive rebase before anyone else sees it. You get
	both a safety net while you work and a readable history afterwards.
\end{moral}

\section{\texorpdfstring{\cmd{cherry-pick}}{cherry-pick}}\label{sec:git-cherry}

\begin{keys}[0.34]
	\cmd{git cherry-pick <sha>}       & Apply one commit here                        \\
	\cmd{git cherry-pick a1b2..c3d4}  & A range                                      \\
	\cmd{git cherry-pick -n <sha>}    & Apply it but do not commit                   \\
	\cmd{git cherry-pick -{}-abort}   & Undo                                         \\
\end{keys}

\begin{eg}
	The standard use is a hotfix that belongs on two branches: commit it on
	\cmd{main}, then cherry-pick it onto the release branch. As with rebase, the
	result is a \emph{new} commit with a different hash --- the same change,
	recorded twice.
\end{eg}

\section{Stashing}\label{sec:git-stash}

\begin{keys}[0.32]
	\cmd{git stash}                  & Put working-tree changes aside; clean the tree \\
	\cmd{git stash -u}               & Include untracked files                        \\
	\cmd{git stash push -m "wip"}    & \dots{} with a label                           \\
	\cmd{git stash list}             & What is stashed                                \\
	\cmd{git stash pop}              & Reapply the newest and remove it from the list \\
	\cmd{git stash apply stash@\{2\}} & Reapply a particular one, keeping it          \\
	\cmd{git stash drop}             & Discard one                                    \\
	\cmd{git stash show -p}          & The diff it holds                              \\
\end{keys}

\begin{note}
	Stashing is for the interruption: you are half-way through something and must
	switch branches now. It is not storage --- an unlabelled stash from three weeks
	ago is indistinguishable from the other four. If work is worth keeping for more
	than an hour, commit it on a branch instead; branches are free
	(\autoref{sec:git-refs}) and they have names.
\end{note}

\section{Worktrees}\label{sec:git-worktree}

\begin{keys}[0.34]
	\cmd{git worktree add ../hotfix main} & A second working directory, on another branch \\
	\cmd{git worktree list}               & Which ones exist                              \\
	\cmd{git worktree remove ../hotfix}   & Tidy up                                       \\
\end{keys}

\begin{moral}
	A worktree is the better answer to ``I need to look at \cmd{main} while I am in
	the middle of this''. It gives you a second directory sharing one object
	database --- no stash, no switching, and your build outputs stay where they
	are. For a project with a long compile, this alone is worth the chapter.
\end{moral}
